% BHU PYQ -- Manually transcribed & error-corrected
\documentclass[11pt,a4paper]{article}

\usepackage[a4paper,top=2.5cm,bottom=2.5cm,left=2.8cm,right=2.8cm,headheight=14pt]{geometry}
\usepackage{amsmath,amssymb}
\usepackage[T1]{fontenc}
\usepackage[utf8]{inputenc}
\usepackage{lmodern}
\usepackage{microtype}
\usepackage{fancyhdr}
\usepackage{enumitem}
\usepackage{hyperref}
\usepackage{lastpage}
\usepackage{parskip}

\hypersetup{colorlinks=true,urlcolor=black,linkcolor=black,
  pdftitle={CHB-502: Inorganic Chemistry-III -- BHU PYQ},pdfauthor={Banaras Hindu University}}

\pagestyle{fancy}\fancyhf{}
\renewcommand{\headrulewidth}{0.4pt}\renewcommand{\footrulewidth}{0.4pt}
\fancyhead[L]{\small   \textit{Banaras Hindu University}}
\fancyhead[R]{\small   \textit{CHB-502: Inorganic Chemistry-III}}
\fancyfoot[C]{\small Page  \thepage\ of \pageref{LastPage}}


\newlist{parts}{enumerate}{2}
\setlist[parts,1]{label=(\alph*),leftmargin=*,itemsep=5pt,topsep=3pt}
\setlist[parts,2]{label=(\roman*),leftmargin=*,itemsep=3pt,topsep=2pt}

\newcommand{\pts}[1]{\hfill   \textit{[#1]}}


\begin{document}


\begin{center}
  {\large\bfseries BANARAS HINDU UNIVERSITY}\\[2pt]
  {\normalsize B.Sc. (Hons.) Semester V Examination, 2013-14}\\[6pt]
  \rule{0.8\linewidth}{0.5pt}\\[6pt]
  {\large\bfseries Paper No. CHB-502: Inorganic Chemistry-III}\\[4pt]
  {\normalsize Time: 3 Hours\hfill Maximum Marks: 70}
\end{center}
\rule{\linewidth}{0.4pt}
\small



\noindent   \textit{Instructions: Answer five questions in total, including Question~1 which is compulsory.}

\normalsize
\bigskip




\noindent   \textbf{Question 1.} Answer all parts of the following: \pts{5 $ \times$ 2 = 10}

\begin{parts}
  \item Why do $3 \text{d}$-elements form stable complexes with nitrogen, oxygen, and fluorine donor ligands, whereas $4 \text{d}$ and $5 \text{d}$ elements prefer phosphorus, sulfur, and chlorine donors?
  \item What is the chelate effect? Give one example.
  \item Why is $[ \text{Co(NH}_3)_6]^{2+}$ oxidized easily to $[ \text{Co(NH}_3)_6]^{3+}$?
  \item Why do $3 \text{d}$ metals form both low-spin and high-spin complexes, whereas $4 \text{d}$ and $5 \text{d}$ metals form only low-spin complexes?
  \item Draw the structure of Zeise's salt.
\end{parts}

\medskip\rule{\linewidth}{0.2pt}




\noindent   \textbf{Question 2.}

\begin{parts}
  \item For the formation of $[ \text{Fe(bipy)}_3]^{2+}$ from $[ \text{Fe(H}_2 \text{O)}_6]^{2+}$ and bipyridyl, the stepwise formation constants lie in the order $K_1 > K_2 < K_3$. Explain this anomaly. \pts{6}
  \item Define and discuss the Jahn-Teller effect. Predict which of the following octahedral complexes show Jahn-Teller distortion: $[ \text{Cr(H}_2 \text{O)}_6]^{3+}$, $[ \text{Cu(H}_2 \text{O)}_6]^{2+}$, and $[ \text{Ni(H}_2 \text{O)}_6]^{2+}$. \pts{8}
\end{parts}

\medskip\rule{\linewidth}{0.2pt}




\noindent   \textbf{Question 3.}

\begin{parts}
  \item Discuss the Crystal Field Theory. Describe the splitting of d-orbitals in octahedral, tetrahedral, and square-planar complexes. What is the relation between octahedral and tetrahedral splitting energy ($\Delta_t = \frac{4}{9}\Delta_o$)? \pts{8}
  \item Comment on the hydration energies of divalent transition metal ions ($M^{2+}$) of the $3 \text{d}$ series. \pts{6}
\end{parts}

\medskip\rule{\linewidth}{0.2pt}




\noindent   \textbf{Question 4.}

\begin{parts}
  \item The reaction between $[ \text{PtCl(dien)}]^+$ and $ \text{I}^-$ is first-order in each reactant. Discuss the mechanism of substitution in square-planar complexes, highlighting the pathway and intermediate involved. \pts{8}
  \item Discuss the factors responsible for the catalytic activity of Wilkinson's catalyst. Describe the mechanism of the hydrogenation of olefins using this catalyst. \pts{6}
\end{parts}

\medskip\rule{\linewidth}{0.2pt}




\noindent   \textbf{Question 5.}

\begin{parts}
  \item What is haptacity? Show that the haptacity of a ligand varies from one organometallic compound to another using cyclopentadienyl and allyl complexes. \pts{8}
  \item Describe the synthesis and mechanism of action of the Ziegler-Natta catalyst for the stereospecific polymerization of propylene. \pts{6}
\end{parts}

\vfill
\rule{\linewidth}{0.4pt}

\begin{center}\small
  \href{https://bhu-exam.vercel.app/}{Click here to visit our website}\\[4pt]
  \href{https://me-aryan.vercel.app/}{Click here to visit our contributor}\\[4pt]
  \href{https://quashan-me.vercel.app/}{Click here to visit our contributor}
\end{center}

\end{document}
